\tcbset{
  appendixexample/.style={
    enhanced jigsaw,
    breakable,
    width=\linewidth,
    colback=black!4,
    colframe=black!60,
    colbacktitle=black!65,
    coltitle=white,
    fonttitle=\small\rmfamily,
    boxrule=0.8pt,
    arc=2pt,
    left=8pt,
    right=8pt,
    top=6pt,
    bottom=6pt,
    before skip=0.45em,
    after skip=0.55em,
  }
}

\subsection{Prompt Template of Initial Planner}
\label{app:planner-initial-prompt}

\begin{tcolorbox}[appendixexample,title={Prompt Template of Initial Planner: $P_{\mathrm{plan}}^{(0)}$}]
\begin{Verbatim}[fontsize=\scriptsize,formatcom=\linespread{0.82}\selectfont,breaklines=true,breakanywhere=true]
You are a memory retrieval planner. Given a user query, decompose it into
targeted retrieval cues for searching the user's conversation history.

## Output format (XML only)

<cues>
  <cue>
    <text>retrieval query text</text>
  </cue>
</cues>

---

## Query

{query}

---

## Rules

Cue text:
- Decompose the query into one or more focused retrieval cues.
- Each cue should target a distinct aspect or entity of the query.
- Always generate at least one `<cue>`. If not sure what to generate,
  copy the original user query.
- Preserve named entities, dates, locations, event types, quantities, and
  temporal semantics.
- Preserve the user's intent as closely as possible. Do not broaden a
  specific question into a generic summary query.
- For questions asking "when", keep the cue as a "when" query about the
  same event.
- For dependent temporal questions, emit the original standalone cue plus
  support cues for the anchor event and target fact.
- For count, comparison, ordering, or duration questions, emit support cues
  that retrieve the underlying items, quantities, event times, or start/end
  anchors.
- Keep the original standalone intent as the first cue, then add supporting
  cues.
- Do not annotate temporal phrases. Temporal extraction is handled by
  deterministic preprocessing after cue generation.

Return XML only, no markdown fences, no text outside the tags.
\end{Verbatim}
\end{tcolorbox}
\captionof{figure}{Prompt template used by the initial planner to decompose the user query into targeted retrieval cues.}

\subsection{Prompt Template of Planner Retry}
\label{app:planner-retry-prompt}

\begin{tcolorbox}[appendixexample,title={Prompt Template of Planner Retry: $P_{\mathrm{plan}}^{(\mathrm{retry})}$}]
\begin{Verbatim}[fontsize=\scriptsize,formatcom=\linespread{0.82}\selectfont,breaklines=true,breakanywhere=true]
You are a memory retrieval planner. A previous retrieval round was
attempted but the synthesizer judged the evidence insufficient.
Your task is to generate new targeted cues that address the identified gap.

## Output format (XML only)

<cues>
  <cue>
    <text>retrieval query text</text>
  </cue>
</cues>

---

## Original query

{query}

## Previous cues

{previous_cues}

## Temporal Evidence Pool
{pool}

## Synthesizer reasoning from previous round

{previous_reason}

---

## Rules

Cue text:
- Generate 1-3 cues that target the missing evidence gap named in the
  synthesizer reasoning and its extracted temporal evidences.
- If reasoning found an anchor/bridge entity/place/time, reuse that entity to
  guide the new cue generation.
- Combine the anchor/bridge with the missing target question relation, e.g.
  added beside, requested from, mailed for, set up through, etc.
- Do not re-search only for the anchor/source if it was already found in reason;
  search for the downstream action/object involving the anchor.
- Turn "no evidence of X near/from/through Y" into a cue for "X
  near/from/through Y". It is your job to recover missing relations through new
  cues.
- Do not repeat previous cues unless the reasoning explicitly asks for
  the same topic with a different temporal anchor or framing.
- Preserve the user's intent as closely as possible. Do not broaden a
  specific question into a generic summary query.
- Keep cues short and search-like; avoid meta questions such as
  "is this linked" or "is this the same item".
- For questions asking "when", keep the cue as a "when" query about the
  same event.
- For count, comparison, ordering, or duration questions, emit support cues
  that retrieve the underlying items, quantities, event times, or start/end
  anchors.
- Do not annotate temporal phrases. Temporal extraction is handled by
  deterministic preprocessing after cue generation.

Return XML only, no markdown fences, no text outside the tags.
\end{Verbatim}
\end{tcolorbox}
\captionof{figure}{Prompt template used by the retry planner to generate new cues from the previous Synthesizer reasoning.}

\subsection{Prompt Template of Synthesizer}
\label{app:synthesizer-prompt}

\begin{tcolorbox}[appendixexample,title={Prompt Template of Synthesizer: $P_{\mathrm{synth}}$}]
\begin{Verbatim}[fontsize=\scriptsize,formatcom=\linespread{0.82}\selectfont,breaklines=true,breakanywhere=true]
You are a memory evidence synthesizer. Given a user query, temporal hints,
previously extracted events, and retrieved conversation sessions, extract
structured evidence and judge whether retrieval is sufficient to answer the
query.

## Output format (XML only, this exact order)

<reasoning>
Concise verdict justification (<=200 words): what decisive evidence is
present, what is missing if any, and why the verdict follows.

</reasoning>

<events>
  <event
    session_id="source session ID; exactly one session per event"
    session_time="YYYY-MM-DDTHH:MM:SS.sssZ"
    event_time="YYYY-MM-DD or 'unknown'">
    A single self-contained evidence statement. Free text, one paragraph,
    no nested tags, no markdown.
  </event>
</events>

<verdict>sufficient|insufficient</verdict>

---

## Normalized Query

{query}

## Retrieved sessions

{sessions}

Retrieved session text preserves original temporal phrases and may annotate
them with parenthesized normalized times, such as `last night (during
2023-05-27 18:00:00 to 2023-05-28 05:59:59)`. These parenthesized values are
deterministic parser aids for interpreting time; they are not additional facts
or evidence beyond the session text.

## Temporal Evidence Pool

{previous_events}

Temporal Evidence Pool are evidence already extracted in earlier retrieval rounds.
They may summarize facts that also appear in the current retrieved sessions.
Use them as context for sufficiency and deduplication, but do not treat them
as new evidence or a substitute for facts grounded in retrieved sessions.
---

## Rules

`Verdict`:
`insufficient` - critical evidence clearly absent; loop continues.
`sufficient` - events reliably cover what the query needs; loop ends.

`event_time`: resolve relative expressions using session timestamps and
temporal hints. Each temporal hint contains the original query phrase and an
absolute UTC timestamp range in `iso_span`. If no explicit date is
recoverable, fall back to the session date. Use `unknown` only if even the
session date is unavailable.

Output rules:
- The final `<events>` output is cumulative. Include all useful
  `<previous_events>` first, unchanged when still valid, then append newly
  extracted events from the current retrieved sessions. Do not delete or
  rewrite previous events unless you find an error that can be corrected by
  current sessions.
- One `<event>` per distinct evidence item, grounded in exactly one session
  (`session_id` must be a single session ID from the input; one session may
  yield multiple events). Always emit partial events even when the verdict is
  `insufficient`. Use `<events/>` only if no useful evidence exists at all.
- If a fact is already covered by previous events and the current sessions add
  no new information, do not repeat it as a new event.
- Preserve temporal relations (before/after/during, duration anchors).
- Event text: free text, one paragraph, no nested tags, no markdown.
- Event text should avoid including temporal phrases when possible; represent
  time in the `event_time` attribute instead.
- Return XML only, no text outside the tags.
- All three attributes (`session_id`, `session_time`, `event_time`) are
  required on every `<event>`.
- For quantity, state, or collection queries, retain all update events
  in chronological order; do not drop later increments or corrections.

\end{Verbatim}
\end{tcolorbox}
\captionof{figure}{Prompt template used by the Synthesizer to extract temporal evidence and decide whether retrieval is sufficient.}

\subsection{Case Study: Temporal Bridge Retrieval}
\label{app:case-study-temporal-bridge}

\begin{tcolorbox}[appendixexample,title={Case Study: Temporal Bridge Retrieval}]
\textbf{Case overview.}
This case illustrates how \sys recovers an answer when the decisive temporal
anchor is not stated in the original query but must be inferred from a bridge
session retrieved in an earlier round.

\smallskip
\textbf{Query.}
What did Haruki read 4 months after sorting through old notebooks and setting
aside pages that could become personal essays?

\smallskip
\textbf{Round 1.}
The initial planner issues cues for the original question, the notebook-sorting
event, and Haruki's reading activities. Retrieval returns several reading
distractors from February, April, May, September, November, and December,
together with the bridge session \texttt{S\_E001}. The key bridge evidence is:
\emph{Haruki: I spent most of the afternoon on the living room floor with old
notebooks and binders, and 2 months ago I made a separate pile of pages that
could become personal essays instead of project records.}

\smallskip
Since \texttt{S\_E001} is timestamped 2024-04-18, the normalized phrase ``2
months ago'' grounds the sorting event in February 2024. The Synthesizer then
infers that four months after February 2024 is June 2024. However, because the
retrieved sessions contain no June reading evidence, it returns
\texttt{insufficient}.

\smallskip
\textbf{Retry cues.}
\begin{Verbatim}[fontsize=\scriptsize,formatcom=\linespread{0.82}\selectfont,breaklines=true,breakanywhere=true]
What did Haruki read in June 2024 (during 2024-06)?
What books did Haruki read between May 2024 and September 2024?
What reading activities did Haruki have in June 2024 (during 2024-06)?
\end{Verbatim}

\smallskip
\textbf{Round 2.}
The second retrieval round finds the target session
\texttt{S\_H2\_E001\_01}, timestamped 2024-06-08:
\emph{Haruki: I read The Death and Life of Great American Cities this month,
and it got me thinking about how to turn my old engineering observations into
personal essays instead of reports.}

\smallskip
\textbf{Final synthesis.}
The final Synthesizer output carries forward the bridge event, appends the June
reading event, and marks the evidence as \texttt{sufficient}. The final answer
is \emph{The Death and Life of Great American Cities}, matching the ground-truth
label. This example highlights the role of Synthesizer reasoning as the
information bridge between rounds: it converts first-round bridge evidence into
a concrete temporal cue that the retry planner can use to retrieve the otherwise
hidden answer session.
\end{tcolorbox}
\captionof{figure}{Case study showing how first-round bridge evidence is transformed into a second-round temporal retrieval cue.}
